
\chapter{Appendixes}
\label{chapter:appendixes}

\newpage

\section{Additional Tables}

Table \ref{tab:emotion_classes} presents the mapping between class numbers, emotion categories, and their corresponding abbreviations used throughout this work. This reference is provided to facilitate the interpretation of the confusion matrices and classification results reported in the \nameref{chapter:results}.

\begin{table}[h]
\centering
\caption{Emotion Classes and Corresponding Labels}
\label{tab:emotion_classes}
\begin{tabular}{c l c}
\hline
\textbf{Class Number} & \textbf{Emotion} & \textbf{Abbreviation} \\
\hline
1 & Angry    & AN \\
2 & Disgust  & DI \\
3 & Fear     & FE \\
4 & Happy    & HA \\
5 & Neutral  & NE \\
6 & Sad      & SA \\
7 & Surprise & SU \\
\hline
\end{tabular}
\end{table}

Table \ref{tab:training_hyperparameters} shows the fixed hyperparameters used for all the experiments. 

\begin{table}[h]
\centering
\caption{Training Hyperparameters and Experimental Configuration}
\label{tab:training_hyperparameters}
\begin{tabular}{lll}
\hline
\textbf{Category} & \textbf{Parameter} & \textbf{Value} \\
\hline

\textbf{Input} & Image resolution & $224 \times 224$ \\
 & Normalization & ImageNet mean/std \\
 & Augmentation & Flip and Random Erasing \\
Model & Architecture & POSTER \\
 & Model sizes & Small / Base / Large \\
 & Head type & Simple / Enhanced \\
\textbf{Optimizer} & Optimizer & Adam + SAM \\
 & Learning rate & $4 \times 10^{-5}$ \\
 & SAM $\rho$ & 0.05 \\
\textbf{Loss} & Loss function & $2\cdot$LSCE + CE \\
 & Label smoothing & 0.2 \\
\textbf{Training} & Epochs & 250 \\
 & Train batch size & 200 \\
 & Validation batch size & 32 \\
\textbf{Scheduler} & LR scheduler & ExponentialLR ($\gamma=0.98$) \\
\textbf{Reproducibility} & Random seed & 123 \\
\hline
\end{tabular}
\end{table}

The computational resources used for all experiments, including both local machines and cloud environments, are summarized in Table~\ref{tab:hardware_used}.

\begin{table}[h]
\centering
\caption{Computational resources used in the experimental setup}
\label{tab:hardware_used}
\renewcommand{\arraystretch}{1.4}
\begin{tabular}{@{}l l l@{}}
\toprule
\textbf{Environment} & \textbf{Component} & \textbf{Specs} \\
\midrule
\multirow{2}{*}{Local}
 & GPU Setup & 2 × NVIDIA GeForce RTX 3060 \\
 & Machine Type & Desktop workstation \\
\midrule
\multirow{2}{*}{Cloud}
 & Platform & Google Colab Pro+ \\
 & GPU & NVIDIA GPU (variable allocation) \\
\bottomrule
\end{tabular}
\end{table}


\section{Source Code}
\label{deliverable:source-code}

%\subsection{General details}
%Pointer or description of the source code can be given in this section..

The complete implementation developed for this thesis, including the modified POSTER architecture, training scripts, and evaluation utilities, is available in a public GitHub repository. This repository \footnote{Source code repository: \url{https://github.com/asvaaron/POSTER}} contains all source code, configuration files, and experiment logs used throughout the research process.

%\section{Gathered Data from the Experiments}
%
%If required, raw data could be presented in this section.
%


\clearpage
\section*{AI Use Statement}
\addcontentsline{toc}{section}{AI Use Statement}

\textbf{AI Use Statement}
Generative AI tools were used only for limited support tasks, including English proofreading/rephrasing and generic programming assistance. AI tools were not used to generate or alter results, datasets, figures, or conclusions. All outputs suggested by AI were reviewed, tested when applicable, and edited by the author, who assumes full responsibility for the content and integrity of this thesis.
